\documentclass[11pt]{article}
\usepackage[utf8]{inputenc}
\usepackage[spanish]{babel}
\usepackage{geometry}
\usepackage{amsmath}
\usepackage{booktabs}
\usepackage{xcolor}
\usepackage{mdframed}

\geometry{a4paper, margin=1in}

\title{Especificaciones Técnicas: HackRF One}
\author{Ingeniería de Radiofrecuencia}
\date{\today}

\newmdenv[
  backgroundcolor=red!8,
  linecolor=red!60,
  linewidth=1pt,
  roundcorner=4pt,
  innertopmargin=6pt,
  innerbottommargin=6pt
]{advertencia}

\newmdenv[
  backgroundcolor=yellow!10,
  linecolor=yellow!60!black,
  linewidth=1pt,
  roundcorner=4pt,
  innertopmargin=6pt,
  innerbottommargin=6pt
]{pendiente}

\begin{document}

\maketitle
\tableofcontents
\newpage

% =========================================================
\section{Ganancias}
% =========================================================

Controladas por separado mediante \texttt{libhackrf} (herramientas como \texttt{hackrf\_transfer} o GNU Radio):

\begin{itemize}
    \item \textbf{LNA (IF RX):} \texttt{hackrf\_set\_lna\_gain()} — 0 a 40 dB, ajustable en pasos de 8 dB. Afecta la señal en la etapa de frecuencia intermedia (IF).
    \item \textbf{VGA (Baseband RX):} \texttt{hackrf\_set\_vga\_gain()} — 0 a 62 dB, ajustable en pasos de 2 dB. Amplificador de ganancia variable en banda base.
    \item \textbf{TX VGA:} \texttt{hackrf\_set\_txvga\_gain()} — 0 a 47 dB, ajustable en pasos de 1 dB.
    \item \textbf{AMP (Front-end RF):} \texttt{hackrf\_set\_amp\_enable()} — 0 (apagado) o 1 (encendido). Activa el amplificador RF de front-end, proporcionando $\approx$+14 dB fijos.
    \begin{pendiente}
    El part number del amplificador varía según la revisión de hardware. El MGA-81563 aparece en esquemáticos de terceros para revisiones tempranas, pero no está confirmado como especificación oficial en la documentación de Great Scott Gadgets.
    \end{pendiente}
    Aplica a la etapa final de TX o etapa inicial de RX.
\end{itemize}

\textbf{Nota:} Dado que el ADC es de 8 bits (rango dinámico limitado a $\approx$48 dB), es crucial gestionar la ganancia correctamente para evitar subir el piso de ruido o causar recorte (clipping) en la conversión analógico-digital.

% =========================================================
\section{Filtros Hardware y Procesamiento}
% =========================================================

\begin{itemize}
    \item \textbf{Filtro Analógico (Baseband LPF):} \texttt{hackrf\_set\_baseband\_filter\_bandwidth()}\\
    Umbrales seleccionables en el transceptor MAX2837 (r1--r8) / MAX2839 (r9): 1.75, 2.5, 3.5, 5, 5.5, 14, 20, 24 y 28 MHz. Ayuda a evitar el aliasing analógico.
    \item \textbf{CPLD y Microcontrolador (Sin DSP interno):}\\
    Utiliza un CPLD Xilinx XC2C64A para enrutamiento de señales y un MCU ARM Cortex-M4 (NXP LPC4320/30). A diferencia de otras SDRs de gama alta, el HackRF One \textbf{no realiza procesamiento de señales digitales interno} (como DDC/DUC o decimación en hardware). Todo el procesamiento matemático recae sobre la CPU del host vía USB.
\end{itemize}

% =========================================================
\section{Rangos de Frecuencia}
% =========================================================

\begin{itemize}
    \item \textbf{Rango Nominal:} 1 MHz a 6 GHz.
    \item \textbf{Arquitectura de Sintetización:}
    \begin{itemize}
        \item 2.3 GHz a 2.7 GHz: Conversión directa utilizando el transceptor MAX2837/MAX2839.
        \item Fuera de esa banda: Utiliza el mezclador RFFC5071 para hacer up-conversion o down-conversion hacia la banda del transceptor.
    \end{itemize}
    \item \textbf{Nota de rendimiento:} Por debajo de 10 MHz, la eficiencia de los transformadores RF (baluns) disminuye.
\end{itemize}

% =========================================================
\section{Sample Rate y Resolución}
% =========================================================

Limitado por la capacidad del ADC/DAC (MAX5864) y el bus USB 2.0.

\begin{table}[h]
    \centering
    \begin{tabular}{llll}
        \toprule
        \textbf{Modo} & \textbf{Resolución I/Q} & \textbf{Sample Rate} & \textbf{Ancho de Banda Máximo} \\
        \midrule
        Estándar  & 8-bit I / 8-bit Q & 2 a 20 MS/s & 20 MHz (Nyquist teórico) \\
        \bottomrule
    \end{tabular}
\end{table}

\textbf{Nota:} Para evitar descartes de paquetes (dropped samples) a 20 MS/s, se recomienda conectar el HackRF a un controlador USB dedicado en el host que no comparta ancho de banda con otros periféricos de alta demanda. Tasas inferiores a 2 MS/s no están soportadas oficialmente.

% =========================================================
\section{Correcciones IQ \& DC}
% =========================================================

\begin{itemize}
    \item \textbf{Pico DC (DC Offset / LO Leakage):} Debido a su arquitectura homodina (Zero-IF / Direct Conversion), el HackRF One presenta una espiga visible en el centro del espectro (0 Hz en banda base). \textbf{Debe corregirse en software} (ej. \textit{DC Blocker} en GNU Radio) o utilizando \textit{Offset Tuning} (sintonizar unos MHz al lado y desplazar digitalmente).
    \item \textbf{Desbalance I/Q:}
    \begin{pendiente}
    Se puede percibir como frecuencias espejo fantasmas al trabajar con señales fuertes. Se recomienda calibrar la ganancia y fase IQ dentro del software host (SDR\#, GNU Radio) ya que el dispositivo no aplica correcciones automáticas de fábrica.
    \end{pendiente}
\end{itemize}

% =========================================================
\section{Conectores y Entradas}
% =========================================================

\begin{table}[h]
    \centering
    \small
    \begin{tabular}{p{0.18\textwidth}p{0.26\textwidth}p{0.44\textwidth}}
        \toprule
        \textbf{Conector} & \textbf{Función} & \textbf{Especificación} \\
        \midrule
        Micro-USB   & Alimentación y datos  & USB 2.0 High Speed (480 Mbps). 5 V / $\approx$500 mA. \\
        ANT (SMA)   & TX/RX principal       & Hembra, 50 $\Omega$. Half-duplex. 1 MHz -- 6 GHz. \\
        CLKIN (SMA) & Referencia externa    & 10 MHz. Onda cuadrada, \textbf{0 V a 3.0 V}. No exceder 3.0 V ni bajar de 0 V. No usar otras frecuencias salvo modificación de firmware. \\
        CLKOUT (SMA)& Reloj de salida       & 10 MHz, 3.3 V. Habilitado con \texttt{hackrf\_clock -o 1}. Para daisy chain entre unidades. \\
        P20, P22, P28 & Headers de expansión & Pines del MCU LPC4320. 3.3 V lógicos, para módulos como PortaPack. \\
        \bottomrule
    \end{tabular}
\end{table}

% =========================================================
\section{Límites de Operación Seguros}
% =========================================================

\begin{advertencia}
\textbf{Advertencias críticas — Susceptibilidad del LNA / Frontend}
\end{advertencia}

\vspace{4pt}
\begin{itemize}
    \item \textbf{Potencia Máxima en RX:} NUNCA exceder \textbf{-5 dBm} ($\approx$0.3 mW) en el puerto de antena. Superarlo destruirá el amplificador de front-end. En teoría, con el AMP desactivado puede tolerar hasta +10 dBm, pero un error de software que lo reactive causaría daño permanente — usar siempre atenuador externo.
    \item \textbf{Potencia TX por banda} (referencia para dimensionar loopback y amplificadores):
    \begin{itemize}
        \item 1--10 MHz: 5 a 15 dBm (aumenta con la frecuencia).
        \item 10--2170 MHz: 5 a 15 dBm (disminuye con la frecuencia).
        \item 2170--2740 MHz: 13 a 15 dBm (mejor zona de rendimiento).
        \item 2740--4000 MHz: 0 a 5 dBm.
        \item 4000--6000 MHz: $-$10 a 0 dBm.
    \end{itemize}
    \textbf{En loopback TX$\to$RX:} usar atenuación suficiente según la banda. Nunca conectar TX directo a RX.
    \item \textbf{Puerto RF sin carga:} No transmitir sin antena o carga ficticia de 50 $\Omega$. La potencia reflejada puede quemar la etapa de TX.
    \item \textbf{Voltajes de Reloj:} Las señales en CLKIN no deben ser negativas ni exceder \textbf{3.0 V}. Una senoidal centrada en 0 V requiere offset DC para cumplir este requisito.
    \item \textbf{Alimentación Antena (Bias-T):} Puede proveer 3.3 V a 50 mA en el conector SMA para alimentar LNAs externos. No activar si el dispositivo conectado presenta cortocircuito en DC.
\end{itemize}

% =========================================================
\section{Clock y Sincronización}
% =========================================================

\begin{itemize}
    \item \textbf{Oscilador Interno:} Cristal de 10 MHz estándar. La estabilidad en ppm no está especificada oficialmente por Great Scott Gadgets; estimaciones comunitarias sugieren 20--30 ppm, susceptible a deriva térmica.
    \item \textbf{Módulo TCXO (Externo/Adicional):} Para uso serio en frecuencias altas o modulación estrecha, acoplar un módulo TCXO de 0.5 ppm en el header P22 o usar CLKIN con referencia externa estable.
    \item \textbf{CLKIN:}
    \begin{itemize}
        \item Señal de 10 MHz, onda cuadrada de 0 a 3.0 V. Alta impedancia.
        \item El dispositivo conmuta automáticamente al reloj externo al iniciar la siguiente operación TX/RX. Verificar detección con \texttt{hackrf\_debug --si5351c -n 0 -r} (salida \texttt{0x01} = clock detectado).
        \item Se puede encadenar directamente CLKOUT de un HackRF One al CLKIN de otro.
    \end{itemize}
    \item \textbf{CLKOUT:} Activar con \texttt{hackrf\_clock -o 1}. Permite sincronizar arrays de múltiples HackRF One con un único reloj maestro.
\end{itemize}

% =========================================================
\section{LEDs Indicadores}
% =========================================================

Seis LEDs montados a lo largo de la placa. Al conectar a USB, \textbf{cuatro deben encenderse}: 3V3, 1V8, RF y USB.

\begin{table}[h]
    \centering
    \begin{tabular}{p{0.10\textwidth}p{0.38\textwidth}p{0.40\textwidth}}
        \toprule
        \textbf{LED} & \textbf{Función} & \textbf{Estado e interpretación} \\
        \midrule
        \textbf{3V3} & Alimentación principal & Encendido al conectar USB. Indica que el regulador de 3.3 V funciona. Si no enciende: fallo de hardware. \\
        \addlinespace
        \textbf{1V8} & Firmware activo + alimentación secundaria & Encendido al conectar USB. Indica que el firmware está corriendo y ha activado fuentes de alimentación internas adicionales (1.8 V). Si no enciende con 3V3 encendido: problema de firmware. \\
        \addlinespace
        \textbf{RF}  & Cadena RF energizada & Encendido al conectar USB. Indica que el firmware ha activado la alimentación del bloque RF. \\
        \addlinespace
        \textbf{USB} & Comunicación USB activa & Encendido cuando el host reconoce el dispositivo correctamente. \\
        \addlinespace
        \textbf{RX}  & Recepción en curso & Encendido \textbf{solo} durante capturas activas. \\
        \addlinespace
        \textbf{TX}  & Transmisión en curso & Encendido \textbf{solo} durante emisión activa. \\
        \bottomrule
    \end{tabular}
\end{table}

\textbf{Nota:} Los colores de los LEDs adyacentes son distintos para facilitar la distinción visual, pero los colores en sí no tienen significado funcional.

% =========================================================
\section{Pinout de Expansión (P20 / P22 / P28)}
% =========================================================

Proporciona acceso a GPIOs del microcontrolador (diseñados para interconectar add-ons como el PortaPack). Operan a \textbf{3.3 V}.

\subsection*{P20 y P28}
\begin{itemize}
    \item Contienen pines para ADC interno del LPC, I2C, SPI y alimentación auxiliar.
    \item Son la columna vertebral de la comunicación paralela para enviar muestras IQ hacia una pantalla táctil (PortaPack) en lugar de ir por USB.
\end{itemize}

\subsection*{P22}
\begin{itemize}
    \item Header de reloj y señales de expansión de bus.
    \item Aquí se instalan los módulos \textbf{TCXO de 0.5 ppm}. Al insertarlo, el reloj primario cambia al de este header, estabilizando el PLL general.
    \item También expone CLKIN alternativo (pin 2), activable con \texttt{hackrf\_clock -c p22}.
\end{itemize}

\textbf{Nota de uso:} No exceder los umbrales de corriente del LPC43xx en ningún pin. Excederlos dañará permanentemente la lógica digital central.

\end{document}